\documentclass[a4paper,11pt]{article}

\usepackage[utf8]{inputenc}
\usepackage[T1]{fontenc}
\usepackage[ngerman]{babel}
\usepackage{amsmath,amssymb}
\usepackage{geometry}
\geometry{margin=2.0cm}
\usepackage{enumitem}
\usepackage{fancyhdr}
\usepackage{xcolor}
\usepackage{tikz}
\usepackage{titlesec}

%================= Dark-Mode Farben (Catppuccin Mocha) =================
<<<<<<< HEAD
\definecolor{base}{HTML}{1E1E2E}
\definecolor{fg}{HTML}{CDD6F4}
\definecolor{accent}{HTML}{89B4FA}
\definecolor{accent2}{HTML}{F5C2E7}
\definecolor{green}{HTML}{A6E3A1}
\definecolor{gridcolor}{HTML}{45475A}
\definecolor{boxbg}{HTML}{313244}

\pagecolor{base}

\newcommand{\transp}{^{\mathsf{T}}}
\newcommand{\inv}{^{-1}}
\newcommand{\pinv}{^{+}}
\newcommand{\diag}{\operatorname{diag}}

\newcommand{\schwierig}[1]{\hfill\textnormal{\color{accent2}\itshape (#1)}}

% Karierte Antwortfläche -- #1 = Höhe (z. B. 8cm).
=======
\definecolor{base}{HTML}{1E1E2E}     % Hintergrund
\definecolor{fg}{HTML}{CDD6F4}       % Text (heller Vordergrund)
\definecolor{accent}{HTML}{89B4FA}   % Überschriften / Akzent
\definecolor{accent2}{HTML}{F5C2E7}  % Schwierigkeitsmarkierung
\definecolor{gridcolor}{HTML}{45475A}% Karo-Gitter (dezent)
\pagecolor{base}

\makeatletter
\newenvironment{psmallmatrix}{\left(\begin{smallmatrix}}{\end{smallmatrix}\right)}
\makeatother

\newcommand{\transp}{^{\mathsf{T}}}
\newcommand{\norm}[1]{\left\lVert #1 \right\rVert}
\newcommand{\inner}[2]{\left\langle #1,\, #2 \right\rangle}
\newcommand{\rang}{\operatorname{rang}}
\newcommand{\spur}{\operatorname{spur}}
\newcommand{\diag}{\operatorname{diag}}
\newcommand{\sgn}{\operatorname{sign}}

% Schwierigkeitsmarkierung
\newcommand{\schwierig}[1]{\hfill\textnormal{\color{accent2}\itshape (#1)}}

% Karierte Antwortfläche (Karopapier) -- #1 = Höhe des Bereichs.
% Durchgehendes Karo-Raster, keine einzelnen Schreiblinien, kein \hrulefill.
>>>>>>> fa386f7 (commit)
\newcommand{\karo}[1]{\par\vspace{0.3cm}\noindent
  \begin{tikzpicture}
    \draw[step=5mm, gridcolor, very thin] (0,0) grid (\linewidth,#1);
  \end{tikzpicture}\par\vspace{0.3cm}}

<<<<<<< HEAD
% Beispiel-Box: vorgerechnetes Beispiel auf farbigem Panel
\newcommand{\beispiel}[1]{\par\medskip\noindent
  \begin{tikzpicture}
    \node[fill=boxbg, rounded corners, inner sep=8pt, text width=\dimexpr\linewidth-18pt\relax,
          align=left] {\color{green}\large\bfseries $\blacktriangleright$ Durchgerechnetes Beispiel\par
          \vspace{2pt}\color{fg}\normalsize #1};
  \end{tikzpicture}\par\medskip}

\titleformat{\section}{\color{accent}\normalfont\Large\bfseries}{}{0pt}{}

\pagestyle{fancy}
\fancyhf{}
\lhead{\color{fg}\small SVD, Cholesky \& Pseudoinverse}
\rhead{\color{fg}\small FLA \& Analytische Geometrie}
\cfoot{\color{fg}\thepage}
\renewcommand{\headrule}{{\color{gridcolor}\hrule height 0.4pt}}

=======
\titleformat{\section}{\color{accent}\normalfont\Large\bfseries}{}{0pt}{}

\pagestyle{fancy}\fancyhf{}
\lhead{\color{fg}\small Fortgeschrittene Lineare Algebra und Analytische Geometrie}
\rhead{\color{fg}\small Übungsblatt 05}
\cfoot{\color{fg}\thepage}
\renewcommand{\headrule}{{\color{gridcolor}\hrule height 0.4pt}}
>>>>>>> fa386f7 (commit)
\setlength{\parindent}{0pt}

\begin{document}
\color{fg}
<<<<<<< HEAD

\thispagestyle{fancy}

\begin{center}
  {\color{accent}\LARGE\bfseries Übungsaufgaben: SVD, Cholesky \& Pseudoinverse}\\[0.4em]
  {\large Fortgeschrittene Lineare Algebra und Analytische Geometrie}\\[0.3em]
  {\color{gridcolor}\rule{\linewidth}{0.4pt}}
\end{center}

\vspace{0.4em}

Übungsaufgaben \emph{ausschließlich} zu den drei Themen \textbf{Cholesky-Zerlegung},
\textbf{Singulärwertzerlegung (SVD)} und \textbf{Pseudoinverse (Moore--Penrose)} -- jeweils
$3$ Aufgaben. Vor jedem Themenblock steht ein \textbf{vollständig durchgerechnetes Beispiel}
(grüner Kasten), an dem das Verfahren Schritt für Schritt vorgemacht wird. Danach kommen die
Aufgaben zum selber Rechnen. Lösungen in \texttt{results/aufgaben/loesungen\_05.pdf}.

\medskip
\textbf{Kurzformeln zum Nachschlagen:}
\begin{itemize}[leftmargin=1.4em, itemsep=2pt]
  \item \textbf{Cholesky} ($A=LL\transp$, nur symm.\ \& positiv definit):
        $\ell_{jj}=\sqrt{a_{jj}-\sum_{k<j}\ell_{jk}^2}$,\quad
        $\ell_{ij}=\big(a_{ij}-\sum_{k<j}\ell_{ik}\ell_{jk}\big)/\ell_{jj}$.
  \item \textbf{SVD} ($A=U\Sigma V\transp$): $V=$ EV von $A\transp A$;\ \ $\sigma_i=\sqrt{\lambda_i}$;\ \
        $\vec u_i=\tfrac{1}{\sigma_i}A\vec v_i$;\ \ $\operatorname{Rang}=\#\{\sigma_i>0\}$.
  \item \textbf{Pseudoinverse} $A\pinv$: aus SVD $A\pinv=V\Sigma\pinv U\transp$;\ \ voller Spaltenrang
        $A\pinv=(A\transp A)\inv A\transp$;\ \ voller Zeilenrang $A\pinv=A\transp(AA\transp)\inv$.
\end{itemize}

% ================================================================
% ============================  TEIL A: CHOLESKY  ================
% ================================================================
\clearpage
\section*{Teil A \;--\; Cholesky-Zerlegung}

\beispiel{%
Bestimme die Cholesky-Zerlegung $A=LL\transp$ von
$A=\begin{pmatrix} 4 & 2 & 2\\ 2 & 5 & 3\\ 2 & 3 & 6\end{pmatrix}$
(symmetrisch). Reihenfolge: $\ell_{11},\ell_{21},\ell_{31},\ell_{22},\ell_{32},\ell_{33}$.
\par\smallskip
\textbf{Spalte 1:}\quad
$\ell_{11}=\sqrt{a_{11}}=\sqrt{4}=2$,\quad
$\ell_{21}=\dfrac{a_{21}}{\ell_{11}}=\dfrac{2}{2}=1$,\quad
$\ell_{31}=\dfrac{a_{31}}{\ell_{11}}=\dfrac{2}{2}=1$.
\par\smallskip
\textbf{Spalte 2:}\quad
$\ell_{22}=\sqrt{a_{22}-\ell_{21}^2}=\sqrt{5-1}=2$,\quad
$\ell_{32}=\dfrac{a_{32}-\ell_{31}\ell_{21}}{\ell_{22}}=\dfrac{3-1\cdot1}{2}=1$.
\par\smallskip
\textbf{Spalte 3:}\quad
$\ell_{33}=\sqrt{a_{33}-\ell_{31}^2-\ell_{32}^2}=\sqrt{6-1-1}=2$.
\par\smallskip
\textbf{Ergebnis:}\quad
$L=\begin{pmatrix} 2 & 0 & 0\\ 1 & 2 & 0\\ 1 & 1 & 2\end{pmatrix}$.
\ \emph{Probe} $LL\transp$: Zeile$_3\cdot$Zeile$_3=1+1+4=6=a_{33}$\,\checkmark, alle Einträge stimmen.
Alle $\ell_{jj}>0$ $\Rightarrow$ $A$ ist tatsächlich \textbf{positiv definit}.}

\medskip
\textbf{Wichtig:} Im Nenner steht immer $\ell_{jj}$, \emph{nicht} $a_{jj}$. Wird ein Radikand
$a_{jj}-\sum\ell_{jk}^2\le 0$, so ist $A$ \emph{nicht} positiv definit und es gibt keine
Cholesky-Zerlegung.

% --------------------------------------------------------------
\clearpage
\section*{Aufgabe 1 \;--\; Cholesky $2\times2$ \schwierig{leicht}}

Gegeben sei $A=\begin{pmatrix} 9 & 3\\ 3 & 5\end{pmatrix}$.
\begin{enumerate}[label=\alph*)]
  \item Bestimmen Sie die Cholesky-Zerlegung $A=LL\transp$.
  \item Begründen Sie mit Ihrem Ergebnis, dass $A$ positiv definit ist.
  \item Machen Sie die Probe $LL\transp=A$.
\end{enumerate}

\karo{9cm}

% --------------------------------------------------------------
\clearpage
\section*{Aufgabe 2 \;--\; Cholesky $3\times3$ \schwierig{mittel}}

Gegeben sei die symmetrische Matrix
\[
A=\begin{pmatrix} 1 & 1 & 1\\ 1 & 5 & 5\\ 1 & 5 & 14\end{pmatrix}.
\]
\begin{enumerate}[label=\alph*)]
  \item Bestimmen Sie $L$ mit $A=LL\transp$ (Reihenfolge $\ell_{11},\ell_{21},\ell_{31},\ell_{22},\ell_{32},\ell_{33}$).
  \item Machen Sie die Probe.
\end{enumerate}

\karo{11cm}

% --------------------------------------------------------------
\clearpage
\section*{Aufgabe 3 \;--\; LGS lösen \& Definitheit testen \schwierig{mittel--schwer}}

\begin{enumerate}[label=\alph*)]
  \item Lösen Sie das Gleichungssystem $A\vec x=\vec b$ mit
        \[
        A=\begin{pmatrix} 4 & 2\\ 2 & 10\end{pmatrix},\qquad
        \vec b=\begin{pmatrix} 6\\ 12\end{pmatrix},
        \]
        \textbf{mit Hilfe der Cholesky-Zerlegung}: erst $A=LL\transp$, dann $L\vec y=\vec b$
        (vorwärts) und $L\transp\vec x=\vec y$ (rückwärts).
  \item Prüfen Sie, ob $B=\begin{pmatrix} 2 & 3\\ 3 & 2\end{pmatrix}$ positiv definit ist, indem
        Sie versuchen, eine Cholesky-Zerlegung zu berechnen. Was beobachten Sie?
\end{enumerate}

\karo{11cm}

% ================================================================
% ==============================  TEIL B: SVD  ===================
% ================================================================
\clearpage
\section*{Teil B \;--\; Singulärwertzerlegung (SVD)}

\beispiel{%
Bestimme die volle SVD $A=U\Sigma V\transp$ von
$A=\begin{pmatrix} 2 & 2\\ -1 & 1\end{pmatrix}$.
\par\smallskip
\textbf{1. $A\transp A$ und ihre Eigenwerte:}\quad
$A\transp A=\begin{pmatrix} 2 & -1\\ 2 & 1\end{pmatrix}\!\begin{pmatrix} 2 & 2\\ -1 & 1\end{pmatrix}
=\begin{pmatrix} 5 & 3\\ 3 & 5\end{pmatrix}$.\;
$\det(A\transp A-\lambda E)=0$: $\lambda=5\pm3\Rightarrow\lambda_1=8,\ \lambda_2=2$.
\par\smallskip
\textbf{2. Singulärwerte:}\quad $\sigma_1=\sqrt{8}=2\sqrt2$,\ \ $\sigma_2=\sqrt2$\ (absteigend).
\par\smallskip
\textbf{3. Rechte Singulärvektoren $V$ (normierte EV von $A\transp A$):}\quad
$\lambda_1{=}8:\ \vec v_1=\tfrac{1}{\sqrt2}(1,1)\transp$;\quad
$\lambda_2{=}2:\ \vec v_2=\tfrac{1}{\sqrt2}(1,-1)\transp$.
\par\smallskip
\textbf{4. Linke Singulärvektoren} $\vec u_i=\tfrac1{\sigma_i}A\vec v_i$:\quad
$\vec u_1=\tfrac{1}{2\sqrt2}\!\begin{pmatrix}2&2\\-1&1\end{pmatrix}\!\tfrac{1}{\sqrt2}\!\begin{pmatrix}1\\1\end{pmatrix}
=\tfrac{1}{4}\begin{pmatrix}4\\0\end{pmatrix}=\begin{pmatrix}1\\0\end{pmatrix}$,\quad
$\vec u_2=\tfrac{1}{\sqrt2}\!\begin{pmatrix}2&2\\-1&1\end{pmatrix}\!\tfrac{1}{\sqrt2}\!\begin{pmatrix}1\\-1\end{pmatrix}
=\tfrac{1}{2}\begin{pmatrix}0\\-2\end{pmatrix}=\begin{pmatrix}0\\-1\end{pmatrix}$.
\par\smallskip
\textbf{Ergebnis:}\quad
$U=\begin{pmatrix} 1 & 0\\ 0 & -1\end{pmatrix},\ \
\Sigma=\begin{pmatrix} 2\sqrt2 & 0\\ 0 & \sqrt2\end{pmatrix},\ \
V=\begin{pmatrix} \tfrac{1}{\sqrt2} & \tfrac{1}{\sqrt2}\\[2pt] \tfrac{1}{\sqrt2} & -\tfrac{1}{\sqrt2}\end{pmatrix}.$
\ Rang $=\#\{\sigma_i>0\}=2$.}

\medskip
\textbf{Merke:} $A\transp A$ \emph{quadriert} die Werte, deshalb $\sigma_i=\sqrt{\lambda_i}$.
$U,V$ sind orthogonal ($U\transp U=V\transp V=E$); $\Sigma$ ist gleich groß wie $A$.

% --------------------------------------------------------------
\clearpage
\section*{Aufgabe 4 \;--\; Singulärwerte \& Rang \schwierig{leicht}}

Gegeben sei $A=\begin{pmatrix} 3 & 0\\ 4 & 0\end{pmatrix}$.
\begin{enumerate}[label=\alph*)]
  \item Berechnen Sie $A\transp A$ und daraus die Singulärwerte $\sigma_1\ge\sigma_2$.
  \item Welchen Rang hat $A$?
  \item Bestimmen Sie eine volle SVD $A=U\Sigma V\transp$ (ergänzen Sie $U$ zu einer
        Orthonormalbasis).
\end{enumerate}

\karo{10cm}

% --------------------------------------------------------------
\clearpage
\section*{Aufgabe 5 \;--\; Volle SVD einer $2\times2$-Matrix \schwierig{mittel}}

Bestimmen Sie die volle Singulärwertzerlegung $A=U\Sigma V\transp$ von
\[
A=\begin{pmatrix} 1 & 2\\ 2 & 1\end{pmatrix}.
\]
Geben Sie $U,\Sigma,V$ explizit an und machen Sie an einem Vektor die Probe, dass
$A\vec v_1=\sigma_1\vec u_1$ gilt.

\karo{12cm}

% --------------------------------------------------------------
\clearpage
\section*{Aufgabe 6 \;--\; SVD einer $3\times2$-Matrix \& Rang-1-Näherung \schwierig{schwer}}

Gegeben sei die Matrix
\[
A=\begin{pmatrix} 1 & 1\\ 1 & 0\\ 0 & 1\end{pmatrix}.
\]
\begin{enumerate}[label=\alph*)]
  \item Bestimmen Sie die (reduzierte) SVD: Singulärwerte $\sigma_1,\sigma_2$, die Matrix $V$
        sowie die linken Singulärvektoren $\vec u_1,\vec u_2$. Welchen Rang hat $A$?
  \item Geben Sie die \textbf{beste Rang-1-Näherung} $A_1=\sigma_1\,\vec u_1\vec v_1\transp$ an
        (Eckart--Young).
\end{enumerate}

\karo{12cm}

% ================================================================
% ========================  TEIL C: PSEUDOINVERSE  ===============
% ================================================================
\clearpage
\section*{Teil C \;--\; Pseudoinverse (Moore--Penrose)}

\beispiel{%
Bestimme die Pseudoinverse von
$A=\begin{pmatrix} 1 & 1\\ 1 & 2\\ 1 & 3\end{pmatrix}$
und löse damit das Ausgleichsproblem $A\vec x\approx\vec b$ mit $\vec b=(1,2,2)\transp$.
\par\smallskip
$A$ hat \textbf{vollen Spaltenrang} ($2$ unabh.\ Spalten) $\Rightarrow$ Linksformel
$A\pinv=(A\transp A)\inv A\transp$.
\par\smallskip
\textbf{1.}\ $A\transp A=\begin{pmatrix} 1&1&1\\ 1&2&3\end{pmatrix}\!\begin{pmatrix} 1&1\\ 1&2\\ 1&3\end{pmatrix}
=\begin{pmatrix} 3 & 6\\ 6 & 14\end{pmatrix}$,\quad
$\det=42-36=6$,\quad
$(A\transp A)\inv=\dfrac{1}{6}\begin{pmatrix} 14 & -6\\ -6 & 3\end{pmatrix}$.
\par\smallskip
\textbf{2.}\ $A\pinv=(A\transp A)\inv A\transp
=\dfrac16\begin{pmatrix} 14 & -6\\ -6 & 3\end{pmatrix}\!\begin{pmatrix} 1&1&1\\ 1&2&3\end{pmatrix}
=\dfrac16\begin{pmatrix} 8 & 2 & -4\\ -3 & 0 & 3\end{pmatrix}$.
\par\smallskip
\emph{Probe} $A\pinv A=\tfrac16\!\begin{pmatrix} 8&2&-4\\ -3&0&3\end{pmatrix}\!\begin{pmatrix} 1&1\\ 1&2\\ 1&3\end{pmatrix}
=\tfrac16\!\begin{pmatrix} 6&0\\ 0&6\end{pmatrix}=E$\,\checkmark.
\par\smallskip
\textbf{3.\ Ausgleichslösung:}\quad
$\vec x=A\pinv\vec b=\tfrac16\!\begin{pmatrix} 8&2&-4\\ -3&0&3\end{pmatrix}\!\begin{pmatrix} 1\\ 2\\ 2\end{pmatrix}
=\tfrac16\begin{pmatrix} 4\\ 3\end{pmatrix}=\begin{pmatrix} 2/3\\ 1/2\end{pmatrix}$.
\ Das ist die Gerade mit kleinstem Fehler $\|A\vec x-\vec b\|$.}

\medskip
\textbf{Merke:} voller \emph{Spalten}rang $\to$ $A\pinv=(A\transp A)\inv A\transp$ (es gilt $A\pinv A=E$);
voller \emph{Zeilen}rang $\to$ $A\pinv=A\transp(AA\transp)\inv$ (es gilt $AA\pinv=E$). Quadratisch \&
invertierbar $\Rightarrow$ $A\pinv=A\inv$.

% --------------------------------------------------------------
\clearpage
\section*{Aufgabe 7 \;--\; Pseudoinverse eines Spaltenvektors \schwierig{leicht}}

Gegeben sei $A=\begin{pmatrix} 1\\ 1\\ 1\end{pmatrix}$ (also $3\times1$).
\begin{enumerate}[label=\alph*)]
  \item Bestimmen Sie $A\pinv$ mit der Formel $A\pinv=(A\transp A)\inv A\transp$.
  \item Berechnen Sie $\vec x=A\pinv\vec b$ für $\vec b=(2,4,6)\transp$. Welche bekannte Größe
        kommt dabei heraus?
\end{enumerate}

\karo{9cm}

% --------------------------------------------------------------
\clearpage
\section*{Aufgabe 8 \;--\; Voller Spaltenrang \schwierig{mittel}}

Gegeben sei
\[
A=\begin{pmatrix} 1 & 0\\ 1 & 1\\ 1 & 2\end{pmatrix}.
\]
\begin{enumerate}[label=\alph*)]
  \item Bestimmen Sie $A\pinv=(A\transp A)\inv A\transp$.
  \item Machen Sie die Probe $A\pinv A=E$.
  \item Lösen Sie das Ausgleichsproblem $A\vec x\approx\vec b$ für $\vec b=(1,1,3)\transp$.
\end{enumerate}

\karo{12cm}

% --------------------------------------------------------------
\clearpage
\section*{Aufgabe 9 \;--\; Voller Zeilenrang \& Minimal-Norm-Lösung \schwierig{schwer}}

Gegeben sei das unterbestimmte System mit
\[
A=\begin{pmatrix} 1 & 0 & 1\\ 0 & 1 & 1\end{pmatrix},\qquad
\vec b=\begin{pmatrix} 1\\ 2\end{pmatrix}.
\]
\begin{enumerate}[label=\alph*)]
  \item $A$ hat vollen Zeilenrang. Bestimmen Sie $A\pinv=A\transp(AA\transp)\inv$.
  \item Berechnen Sie die Lösung kleinster Norm $\vec x=A\pinv\vec b$ und prüfen Sie $A\vec x=\vec b$.
\end{enumerate}

\karo{12cm}
=======
\thispagestyle{fancy}

\begin{center}
  {\color{accent}\LARGE\bfseries Übungsblatt 05 -- Komplettabdeckung des Stoffs}\\[0.4em]
  {\large Fortgeschrittene Lineare Algebra und Analytische Geometrie}\\[0.8em]
\end{center}

\noindent
\begin{tabular}{@{}p{0.6\textwidth}p{0.35\textwidth}@{}}
  \textbf{Name:} \textcolor{fg}{\hrulefill} & \textbf{Datum:} \textcolor{fg}{\hrulefill}
\end{tabular}

\vspace{0.4em}{\color{gridcolor}\hrule}\vspace{0.6em}

\textit{Dieses Blatt deckt den \emph{gesamten} Vorlesungsstoff mit \textbf{neuen} Aufgaben ab und
ergänzt die Themen, die auf den Blättern 01--04 dünn waren: symbolisches Rechnen \& Matrix-Rechenregeln,
Zerlegung in symmetrischen/schiefsymmetrischen Anteil, Orthogonalität der vier Unterräume, äußeres
Produkt (Rang-1), Nachweis innerer Produkte. Außerdem: Normen \& Skalarprodukt, LGS-Lösbarkeit (Gauß),
CR-/$PA{=}LU$-/Cholesky-/QR-Zerlegung, Eigenwerte \& Diagonalisierung, Eigenwert-Regeln,
Spektralzerlegung, Definitheit, inneres Produkt $\vec x\transp A\vec y$ \& Mahalanobis, Untervektorräume,
volle \& reduzierte SVD, Pseudoinverse \& Ausgleich, lineare Regression, PCA, Hyperebene \&
Margin-Classifier, Finite Differenzen ($K_n,T_n,B_n,C_n$) und LoRA. Notation wie in der Vorlesung;
$E$ = Einheitsmatrix. „Glatte`` Wurzeln bitte ausrechnen, „krumme`` dürfen als Wurzel/Bruch stehen
bleiben. Begleitend: \texttt{openbook/matrizen.pdf}.}

\vspace{0.6em}

%==================================================================
\clearpage
\section*{Aufgabe 1 -- Symbolisches Rechnen \& Matrix-Rechenregeln\schwierig{mittel}}
$A,B$ seien quadratisch und invertierbar.
\begin{enumerate}[label=\alph*)]
  \item Vereinfachen Sie soweit wie möglich:
        \[ \text{(i)}\ \big((A\transp)^{-1}\big)\transp, \qquad
           \text{(ii)}\ (AB)^{-1}A, \qquad
           \text{(iii)}\ B\,(B^{-1}A B)\,B^{-1}. \]
  \item Welche der folgenden Gleichungen gilt für \emph{alle} quadratischen $A,B$? Begründen Sie kurz
        bzw.\ geben Sie ein Gegenbeispiel.
        \[ \text{(1)}\ AB=BA,\quad \text{(2)}\ (A+B)\transp=A\transp+B\transp,\quad
           \text{(3)}\ (AB)\transp=A\transp B\transp,\quad \text{(4)}\ (A+B)^2=A^2+2AB+B^2. \]
  \item Zeigen Sie: Für jede quadratische Matrix $A$ ist $S=A+A\transp$ symmetrisch und $W=A-A\transp$
        schiefsymmetrisch ($W\transp=-W$). Folgern Sie, dass sich jede quadratische Matrix als Summe einer
        symmetrischen und einer schiefsymmetrischen Matrix schreiben lässt.
  \item Zeigen Sie: Ist $A$ invertierbar und \emph{symmetrisch}, so ist auch $A^{-1}$ symmetrisch.
\end{enumerate}
\karo{8.5cm}

%==================================================================
\clearpage
\section*{Aufgabe 2 -- Quickies (Wahr oder Falsch)\schwierig{leicht}}
Entscheiden Sie für jede Aussage, ob sie \emph{wahr} oder \emph{falsch} ist, und begründen Sie kurz.
\begin{enumerate}[label=\alph*)]
  \item Für quadratische $A,B$ gilt stets $\det(A+B)=\det(A)+\det(B)$.
  \item Ähnliche Matrizen ($B=S^{-1}AS$) besitzen dieselben Eigenwerte.
  \item Die Spalten einer orthogonalen Matrix bilden eine Orthonormalbasis.
  \item Eine $m\times n$-Matrix mit $m<n$ kann vollen Spaltenrang besitzen.
  \item Die Singulärwerte einer Matrix sind ihre Eigenwerte.
  \item Ist $\lambda$ Eigenwert von $A$, so ist $\lambda^2$ Eigenwert von $A^2$.
  \item Der Nullraum $N(A)$ steht orthogonal zum Zeilenraum $C(A\transp)$.
  \item Jede Matrix der Form $A\transp A$ ist symmetrisch und positiv semidefinit.
\end{enumerate}
\karo{7cm}

%==================================================================
\clearpage
\section*{Aufgabe 3 -- Normen, Skalarprodukt \& Winkel\schwierig{leicht}}
Gegeben seien
\[
  \vec x = \begin{pmatrix} 2 \\ -1 \\ 2 \end{pmatrix},\qquad
  \vec y = \begin{pmatrix} 1 \\ 2 \\ 2 \end{pmatrix}.
\]
\begin{enumerate}[label=\alph*)]
  \item Berechnen Sie $\norm{\vec x}_1$, $\norm{\vec x}_2$ und $\norm{\vec x}_\infty$.
  \item Berechnen Sie $\inner{\vec x}{\vec y}=\vec x\transp\vec y$. Sind $\vec x,\vec y$ orthogonal?
  \item Überprüfen Sie die Cauchy-Schwarz-Ungleichung
        $|\inner{\vec x}{\vec y}|\le\norm{\vec x}_2\,\norm{\vec y}_2$.
  \item Bestimmen Sie den Winkel $\alpha$ zwischen $\vec x$ und $\vec y$ (über $\cos\alpha$).
\end{enumerate}
\karo{6.5cm}

%==================================================================
\clearpage
\section*{Aufgabe 4 -- Lösbarkeit von LGS (Gauß, $3\times3$)\schwierig{mittel}}
Untersuchen Sie jedes System mit dem Gauß-Verfahren auf Lösbarkeit und geben Sie die Lösungsmenge an.
\begin{enumerate}[label=\alph*)]
  \item $\begin{aligned}[t] x_1+x_2+x_3&=6\\ x_1+2x_2+3x_3&=14\\ x_1+3x_2+6x_3&=25 \end{aligned}$
  \item $\begin{aligned}[t] x_1+x_2+x_3&=1\\ 2x_1+2x_2+2x_3&=3\\ x_1\phantom{+x_2}+x_3&=0 \end{aligned}$
  \item $\begin{aligned}[t] x_1+x_2+x_3&=3\\ \phantom{x_1+{}}x_2+2x_3&=3\\ x_1+2x_2+3x_3&=6 \end{aligned}$
\end{enumerate}
\karo{8.5cm}

%==================================================================
\clearpage
\section*{Aufgabe 5 -- CR-Zerlegung\schwierig{mittel}}
Gegeben sei
\[
  A = \begin{pmatrix} 1 & 2 & 0 & 1 \\ 2 & 4 & 1 & 4 \\ 1 & 2 & 1 & 3 \end{pmatrix}\in\mathbb{R}^{3\times4}.
\]
\begin{enumerate}[label=\alph*)]
  \item Bringen Sie $A$ auf reduzierte Zeilenstufenform und bestimmen Sie den Rang $r$.
  \item Geben Sie eine Basis des Spaltenraums $C(A)$ an.
  \item Geben Sie die CR-Zerlegung $A=CR$ an und machen Sie die Probe.
\end{enumerate}
\karo{8cm}

%==================================================================
\clearpage
\section*{Aufgabe 6 -- Die vier Unterräume \& ihre Orthogonalität\schwierig{schwer}}
Verwenden Sie wieder $A=\begin{psmallmatrix}1&2&0&1\\2&4&1&4\\1&2&1&3\end{psmallmatrix}$ aus Aufgabe~5 (Rang $r=2$).
\begin{enumerate}[label=\alph*)]
  \item Geben Sie die Dimensionen aller vier Unterräume an
        ($\dim C(A),\ \dim C(A\transp),\ \dim N(A),\ \dim N(A\transp)$).
  \item Bestimmen Sie eine Basis des Nullraums $N(A)$.
  \item Bestätigen Sie an Ihren Basisvektoren, dass $N(A)$ orthogonal zum Zeilenraum $C(A\transp)$ steht.
  \item Bestimmen Sie eine Basis des Linksnullraums $N(A\transp)$
        (Vektor $\vec y$ mit $\vec y\transp A = \vec 0\transp$).
\end{enumerate}
\karo{9cm}

%==================================================================
\clearpage
\section*{Aufgabe 7 -- Äußeres Produkt \& Rang-1-Matrizen\schwierig{mittel}}
Gegeben sei
\[
  A = \begin{pmatrix} 2 & 4 & -2 \\ 3 & 6 & -3 \\ 1 & 2 & -1 \end{pmatrix}.
\]
\begin{enumerate}[label=\alph*)]
  \item Zeigen Sie, dass $A$ als äußeres Produkt $A=\vec u\,\vec v\transp$ zweier Vektoren geschrieben
        werden kann, und geben Sie $\vec u$ und $\vec v$ an.
  \item Welchen Rang hat $A$? Wie groß ist $\det(A)$?
  \item Eine Rang-1-Matrix $\vec u\vec v\transp$ besitzt genau einen von $0$ verschiedenen Eigenwert
        $\lambda=\vec v\transp\vec u$. Bestimmen Sie $\lambda$ und überprüfen Sie $\spur(A)=\lambda$.
\end{enumerate}
\karo{7.5cm}

%==================================================================
\clearpage
\section*{Aufgabe 8 -- Determinante (Sarrus), Spur \& Rechenregeln\schwierig{mittel}}
\begin{enumerate}[label=\alph*)]
  \item Berechnen Sie $\det(A)$ mit der Regel von Sarrus und geben Sie $\spur(A)$ an für
        \[ A = \begin{pmatrix} 2 & -1 & 0 \\ 1 & 3 & 1 \\ 0 & 2 & 2 \end{pmatrix}. \]
  \item Für zwei $3\times3$-Matrizen gelte $\det(A)=10$ und $\det(B)=-2$. Bestimmen Sie
        $\det(AB)$, $\det(2A)$, $\det(A^{-1})$, $\det(A\transp)$ und $\det(B^{-1}A)$.
\end{enumerate}
\karo{7.5cm}

%==================================================================
\clearpage
\section*{Aufgabe 9 -- $2\times2$-Inverse \& Rechenregel-Beweis\schwierig{mittel}}
\begin{enumerate}[label=\alph*)]
  \item Bestimmen Sie mit der $2\times2$-Formel die Inverse von
        $B=\begin{psmallmatrix}3&5\\2&4\end{psmallmatrix}$ und machen Sie die Probe $BB^{-1}=E$.
  \item Beweisen Sie für invertierbare $A,B$ die Regel $(AB)^{-1}=B^{-1}A^{-1}$.
\end{enumerate}
\karo{7cm}

%==================================================================
\clearpage
\section*{Aufgabe 10 -- LU mit Zeilentausch ($PA=LU$)\schwierig{schwer}}
Gegeben seien
\[
  A = \begin{pmatrix} 0 & 1 & 1 \\ 2 & 4 & 6 \\ 1 & 3 & 5 \end{pmatrix},\qquad
  b = \begin{pmatrix} 2 \\ 12 \\ 9 \end{pmatrix}.
\]
\begin{enumerate}[label=\alph*)]
  \item Warum ist das Pivotelement $a_{11}$ ein Problem? Geben Sie die Permutationsmatrix $P$ an.
  \item Bestimmen Sie $L$ und $U$ mit $PA=LU$.
  \item Bestimmen Sie $\det(A)$ (Vorzeichen durch den Zeilentausch beachten!).
  \item Lösen Sie $A\vec x=b$ über $PA\vec x=Pb$ (vorwärts $L\vec y=Pb$, rückwärts $U\vec x=\vec y$).
\end{enumerate}
\karo{9cm}

%==================================================================
\clearpage
\section*{Aufgabe 11 -- Cholesky-Zerlegung\schwierig{mittel}}
Gegeben seien
\[
  A = \begin{pmatrix} 4 & 2 & -2 \\ 2 & 10 & 2 \\ -2 & 2 & 6 \end{pmatrix},\qquad
  b = \begin{pmatrix} 4 \\ 14 \\ 6 \end{pmatrix}.
\]
\begin{enumerate}[label=\alph*)]
  \item Zeigen Sie mit den führenden Hauptminoren, dass $A$ positiv definit ist.
  \item Bestimmen Sie $L$ mit $A=LL\transp$.
  \item Lösen Sie $A\vec x=b$ über $L\vec y=b$ und $L\transp\vec x=\vec y$.
  \item Erklären Sie, warum $M=\begin{psmallmatrix}1&2\\2&1\end{psmallmatrix}$ \emph{keine}
        Cholesky-Zerlegung besitzt.
\end{enumerate}
\karo{8.5cm}

%==================================================================
\clearpage
\section*{Aufgabe 12 -- QR über Gram-Schmidt\schwierig{schwer}}
Gegeben seien
\[
  A = \begin{pmatrix} 1 & 3 \\ 2 & 0 \\ 2 & 3 \end{pmatrix},\qquad b = \begin{pmatrix} 4 \\ 2 \\ 5 \end{pmatrix}.
\]
\begin{enumerate}[label=\alph*)]
  \item Bestimmen Sie mit dem Gram-Schmidt-Verfahren eine \emph{orthonormale} Basis $\vec q_1,\vec q_2$
        der Spalten von $A$ und geben Sie $Q$ an.
  \item Bestimmen Sie $R=Q\transp A$ und machen Sie die Probe $QR=A$.
  \item Lösen Sie $A\vec x=b$ über $R\vec x=Q\transp b$ (Rückwärtseinsetzen).
\end{enumerate}
\karo{8.5cm}

%==================================================================
\clearpage
\section*{Aufgabe 13 -- Inneres Produkt nachweisen\schwierig{mittel}}
Eine bilineare Form $\inner{\vec x}{\vec y}=\vec x\transp A\vec y$ auf $\mathbb{R}^2$ ist genau dann ein
inneres Produkt, wenn $A$ symmetrisch und positiv definit ist. Entscheiden Sie für jede Form, ob sie ein
inneres Produkt ist, und begründen Sie (Symmetrie + Definitheit über die führenden Hauptminoren).
\begin{enumerate}[label=\alph*)]
  \item $\inner{\vec x}{\vec y}=\vec x\transp\begin{psmallmatrix}2&1\\1&1\end{psmallmatrix}\vec y$
  \item $\inner{\vec x}{\vec y}=\vec x\transp\begin{psmallmatrix}1&2\\2&1\end{psmallmatrix}\vec y$
  \item $\inner{\vec x}{\vec y}=x_1y_1-(x_1y_2+x_2y_1)+2x_2y_2$ \\
        (schreiben Sie die Form zuerst als $\vec x\transp A\vec y$).
\end{enumerate}
Geben Sie für die \emph{nicht} positiv definite Form einen Vektor $\vec x\neq\vec 0$ mit
$\vec x\transp A\vec x<0$ an.
\karo{8cm}

%==================================================================
\clearpage
\section*{Aufgabe 14 -- Orthogonale Matrizen\schwierig{mittel}}
Gegeben sei
\[
  Q = \frac{1}{3}\begin{pmatrix} 2 & -2 & 1 \\ 2 & 1 & -2 \\ -1 & -2 & -2 \end{pmatrix}.
\]
\begin{enumerate}[label=\alph*)]
  \item Zeigen Sie, dass $Q$ orthogonal ist (Spalten paarweise orthonormal).
  \item Bestimmen Sie $Q^{-1}$ und $\det(Q)$. Handelt es sich um eine Drehung oder eine Spiegelung?
  \item Überprüfen Sie die Längenerhaltung $\norm{Q\vec x}=\norm{\vec x}$ für $\vec x=(1,2,2)\transp$.
\end{enumerate}
\karo{7.5cm}

%==================================================================
\clearpage
\section*{Aufgabe 15 -- Eigenwerte, Diagonalisierung \& Matrixpotenz\schwierig{mittel}}
Gegeben sei
\[
  A = \begin{pmatrix} 4 & -2 \\ 1 & 1 \end{pmatrix}.
\]
\begin{enumerate}[label=\alph*)]
  \item Bestimmen Sie die Eigenwerte über das charakteristische Polynom $\chi_A(\lambda)$ und zu jedem
        Eigenwert einen Eigenvektor.
  \item Geben Sie $S$ und $D$ mit $A=SDS^{-1}$ an (inklusive $S^{-1}$).
  \item Leiten Sie $A^k=SD^kS^{-1}$ her und berechnen Sie damit $A^3$.
  \item Überprüfen Sie $\spur(A)=\sum_i\lambda_i$ und $\det(A)=\prod_i\lambda_i$.
\end{enumerate}
\karo{8.5cm}

%==================================================================
\clearpage
\section*{Aufgabe 16 -- Eigenwert-Regeln (Dreiecksmatrix)\schwierig{mittel}}
Gegeben sei die obere Dreiecksmatrix
\[
  A = \begin{pmatrix} 2 & 5 & 1 \\ 0 & 3 & 4 \\ 0 & 0 & -1 \end{pmatrix}.
\]
\begin{enumerate}[label=\alph*)]
  \item Geben Sie die Eigenwerte direkt an (Begründung) und bestätigen Sie
        $\spur(A)=\sum_i\lambda_i$ sowie $\det(A)=\prod_i\lambda_i$.
  \item Bestimmen Sie \emph{ohne} neue charakteristische Gleichung die Eigenwerte von $A+3E$, von
        $A^{-1}$ und von $A^2$ (Regeln $\lambda\mapsto\lambda+\gamma,\ \lambda^{-1},\ \lambda^2$).
  \item Ist $A$ invertierbar? Begründen Sie über die Eigenwerte.
\end{enumerate}
\karo{7.5cm}

%==================================================================
\clearpage
\section*{Aufgabe 17 -- Spektralzerlegung ($3\times3$)\schwierig{schwer}}
Gegeben sei die symmetrische Matrix
\[
  A = \begin{pmatrix} 5 & 0 & 0 \\ 0 & 3 & 1 \\ 0 & 1 & 3 \end{pmatrix}.
\]
\begin{enumerate}[label=\alph*)]
  \item Bestimmen Sie die Eigenwerte und eine \emph{orthonormale} Eigenvektorbasis
        (nutzen Sie die Blockstruktur).
  \item Geben Sie die Spektralzerlegung $A=Q\Lambda Q\transp$ an.
  \item Berechnen Sie $A^{-1}=Q\Lambda^{-1}Q\transp$.
\end{enumerate}
\karo{8.5cm}

%==================================================================
\clearpage
\section*{Aufgabe 18 -- Definitheit klassifizieren\schwierig{mittel}}
Klassifizieren Sie die folgenden symmetrischen Matrizen (positiv/negativ (semi-)definit oder indefinit)
jeweils über die führenden Hauptminoren \emph{und} die Eigenwerte:
\[
  A=\begin{psmallmatrix}3&1\\1&3\end{psmallmatrix},\quad
  B=\begin{psmallmatrix}-2&1\\1&-3\end{psmallmatrix},\quad
  C=\begin{psmallmatrix}1&2\\2&4\end{psmallmatrix},\quad
  D=\begin{psmallmatrix}1&3\\3&1\end{psmallmatrix}.
\]
Geben Sie für die \emph{indefinite} Matrix zusätzlich einen Vektor $\vec x$ mit
$\vec x\transp(\text{Matrix})\vec x<0$ an.
\karo{8cm}

%==================================================================
\clearpage
\section*{Aufgabe 19 -- Inneres Produkt $\vec x\transp A\vec y$ \& Mahalanobis\schwierig{mittel}}
\begin{enumerate}[label=\alph*)]
  \item Ist $\inner{\vec x}{\vec y}=\vec x\transp A\vec y$ mit $A=\begin{psmallmatrix}2&-1\\-1&2\end{psmallmatrix}$
        ein inneres Produkt? Begründen Sie (Symmetrie + Definitheit).
  \item Berechnen Sie mit diesem inneren Produkt $\inner{\vec u}{\vec v}$ und die Norm $\norm{\vec u}_A$
        für $\vec u=(1,1)\transp,\ \vec v=(1,-1)\transp$. Sind $\vec u,\vec v$ orthogonal bzgl.\ $A$?
  \item Berechnen Sie den Mahalanobis-Abstand $\sqrt{\vec x\transp\Sigma^{-1}\vec x}$ für
        $\vec x=(4,3)\transp,\ \Sigma=\begin{psmallmatrix}16&0\\0&9\end{psmallmatrix}$.
\end{enumerate}
\karo{8cm}

%==================================================================
\clearpage
\section*{Aufgabe 20 -- Untervektorraum-Test\schwierig{leicht}}
Entscheiden Sie für jede Menge, ob sie ein Untervektorraum ist, und begründen Sie kurz
(Kriterien: $\vec 0\in U$, abgeschlossen unter $+$ und Skalarmultiplikation).
\begin{enumerate}[label=\alph*)]
  \item $U_1=\{(x,y)\transp\in\mathbb{R}^2 : 3x-y=0\}$
  \item $U_2=\{(x,y)\transp\in\mathbb{R}^2 : xy=0\}$
  \item $U_3=\{(x,y,z)\transp\in\mathbb{R}^3 : x-y+2z=0\}$
  \item $U_4=\{(x,y,z)\transp\in\mathbb{R}^3 : x+y+z=1\}$
  \item $U_5=\{(x,y)\transp\in\mathbb{R}^2 : x\ge 0,\ y\ge 0\}$
\end{enumerate}
\karo{7.5cm}

%==================================================================
\clearpage
\section*{Aufgabe 21 -- SVD (volle \& reduzierte Form)\schwierig{schwer}}
Gegeben sei
\[
  A = \begin{pmatrix} 1 & 1 \\ 1 & 1 \\ 1 & -1 \end{pmatrix}.
\]
\begin{enumerate}[label=\alph*)]
  \item Berechnen Sie $A\transp A$ und bestimmen Sie die Singulärwerte (absteigend).
  \item Bestimmen Sie $V$, $\Sigma$ und $U=[\vec u_1\ \vec u_2\ \vec u_3]$ der \emph{vollen} SVD
        $A=U\Sigma V\transp$ (Dimensionen: $U\colon3\times3$, $\Sigma\colon3\times2$, $V\colon2\times2$).
  \item Geben Sie die \emph{reduzierte} SVD an und bestätigen Sie
        $A=\sigma_1\vec u_1\vec v_1\transp+\sigma_2\vec u_2\vec v_2\transp$.
\end{enumerate}
\karo{8.5cm}

%==================================================================
\clearpage
\section*{Aufgabe 22 -- Pseudoinverse \& Ausgleichsrechnung\schwierig{mittel}}
Gegeben seien die Matrix mit vollem Spaltenrang und der Messvektor
\[
  A = \begin{pmatrix} 1 & 0 \\ 1 & 1 \\ 1 & 2 \end{pmatrix},\qquad
  b = \begin{pmatrix} 1 \\ 1 \\ 4 \end{pmatrix}.
\]
\begin{enumerate}[label=\alph*)]
  \item Berechnen Sie $A\transp A$ und $(A\transp A)^{-1}$.
  \item Bestimmen Sie die Moore-Penrose-Pseudoinverse $A^{+}=(A\transp A)^{-1}A\transp$ und
        verifizieren Sie $A^{+}A=E$.
  \item Lösen Sie das überbestimmte System $A\vec x\approx b$ bestmöglich
        ($\vec x=A^{+}b$, kleinste Fehlerquadrate).
\end{enumerate}
\karo{8cm}

%==================================================================
\clearpage
\section*{Aufgabe 23 -- Lineare Regression \& Deutung\schwierig{schwer}}
Zu den Datenpunkten $(x,y)\in\{(0,1),(1,1),(2,4),(3,4)\}$ soll die Ausgleichsgerade
$y=c_0+c_1x$ bestimmt werden.
\begin{enumerate}[label=\alph*)]
  \item Stellen Sie die Designmatrix $X$ und $\vec y$ auf und notieren Sie die Normalengleichung
        $X\transp X\,\vec c=X\transp\vec y$
        (mit $X\transp X=\begin{psmallmatrix}n&\sum x_i\\\sum x_i&\sum x_i^2\end{psmallmatrix}$).
  \item Lösen Sie nach $\vec c$ auf und geben Sie die Geradengleichung an.
  \item \emph{Deutung:} Was bedeutet ein Bestimmtheitsmaß $R^2\approx1$, und was ein kleiner
        $p$-Wert ($<0{,}05$) für die Güte des Modells?
\end{enumerate}
\karo{8.5cm}

%==================================================================
\clearpage
\section*{Aufgabe 24 -- Hauptkomponentenanalyse (PCA)\schwierig{schwer}}
Gegeben seien die Datenpunkte (als Zeilen)
\[
  (3,2),\quad (2,3),\quad (-2,-3),\quad (-3,-2).
\]
\begin{enumerate}[label=\alph*)]
  \item Zeigen Sie, dass die Daten bereits zentriert sind (Spaltenmittel $=0$).
  \item Bestimmen Sie die Kovarianzmatrix $C=\tfrac{1}{n-1}A_c\transp A_c$.
  \item Bestimmen Sie die Eigenwerte von $C$ und die erste Hauptkomponente (normierter Eigenvektor zum
        größten Eigenwert).
  \item Welcher Anteil der Gesamtvarianz wird durch PC1 erklärt?
\end{enumerate}
\karo{8.5cm}

%==================================================================
\clearpage
\section*{Aufgabe 25 -- Hyperebene, Abstand \& Margin-Classifier\schwierig{mittel}}
\textbf{Teil 1.} Gegeben sei die Ebene $E:\ 2x_1-x_2+2x_3=6$.
\begin{enumerate}[label=\alph*)]
  \item Geben Sie einen Normalenvektor $\vec n$ und die Dimension von $E$ als Hyperebene im
        $\mathbb{R}^3$ an. Verläuft $E$ durch den Ursprung?
  \item Liegt $p=(2,1,0)\transp$ auf $E$? Berechnen Sie den Abstand $d$ von $p$ zu $E$.
\end{enumerate}
\textbf{Teil 2.} Zwei Klassen: $\;A:\ (3,1),(3,-1)\qquad B:\ (-1,1),(-1,-1)$.
\begin{enumerate}[label=\alph*),start=3]
  \item Bestimmen Sie $\vec n$ (Richtung zwischen den Klassenmitteln) und $b=-\vec n\transp\vec m$
        ($\vec m=$ Mittelpunkt der Klassenmittel) und geben Sie die trennende Ebene
        $\vec n\transp\vec x+b=0$ an.
  \item Klassifizieren Sie den Testpunkt $p=(0,2)\transp$ über das Vorzeichen von $\vec n\transp p+b$.
  \item Bestimmen Sie den Margin (kleinster Abstand eines Klassenpunktes zur Ebene).
\end{enumerate}
\karo{8.5cm}

%==================================================================
\clearpage
\section*{Aufgabe 26 -- Finite Differenzen \& die vier Matrizen\schwierig{schwer}}
\begin{enumerate}[label=\alph*)]
  \item Geben Sie den zentralen Differenzenquotienten für $u''(x)$ an und leiten Sie für $-u''(x)=f$
        mit $u=0$ am Rand das LGS $T\vec u=h^2\vec f$ her ($-u_{i-1}+2u_i-u_{i+1}=h^2f_i$).
  \item Lösen Sie konkret $-u''(x)=12$, $u(0)=u(1)=0$ mit $x_1=0{,}25,\ x_2=0{,}5,\ x_3=0{,}75$
        ($h=\tfrac14$, also $T=K_3$). Nutzen Sie die Symmetrie $u_1=u_3$.
  \item Klassifizieren Sie die vier $3\times3$-Matrizen
        \[
          K_3=\begin{psmallmatrix}2&-1&0\\-1&2&-1\\0&-1&2\end{psmallmatrix},\
          T_3=\begin{psmallmatrix}1&-1&0\\-1&2&-1\\0&-1&2\end{psmallmatrix},\
          B_3=\begin{psmallmatrix}1&-1&0\\-1&2&-1\\0&-1&1\end{psmallmatrix},\
          C_3=\begin{psmallmatrix}2&-1&-1\\-1&2&-1\\-1&-1&2\end{psmallmatrix}.
        \]
        Welche sind invertierbar/positiv definit, welche singulär? Geben Sie für die singulären je
        einen Nullraumvektor an.
\end{enumerate}
\karo{8.5cm}

%==================================================================
\clearpage
\section*{Aufgabe 27 -- LoRA / beste Rang-$k$-Näherung\schwierig{mittel}}
\textbf{Teil 1.} Gegeben sei $A=\begin{psmallmatrix}5&0\\0&4\\0&0\end{psmallmatrix}$ mit Singulärwerten
$\sigma_1=5$ ($\vec u_1=(1,0,0)\transp,\ \vec v_1=(1,0)\transp$) und
$\sigma_2=4$ ($\vec u_2=(0,1,0)\transp,\ \vec v_2=(0,1)\transp$).
\begin{enumerate}[label=\alph*)]
  \item Bestimmen Sie die beste Rang-1-Näherung $A_1=\sigma_1\vec u_1\vec v_1\transp$.
  \item Wie groß sind die Fehler $\norm{A-A_1}_2$ (Spektralnorm) und $\norm{A-A_1}_F$ (Frobenius)?
\end{enumerate}
\textbf{Teil 2.} Ein Graustufenbild liege als $A\in\mathbb{R}^{5\times4}$ vor mit Singulärwerten
$\sigma_1=12,\ \sigma_2=8,\ \sigma_3=5,\ \sigma_4=2$.
\begin{enumerate}[label=\alph*),start=3]
  \item Welchen Rang hat $A$? Geben Sie die Fehler $\norm{A-A_2}_2$ und $\norm{A-A_2}_F$ der besten
        Rang-2-Näherung $A_2$ an.
  \item Wie viele Zahlen muss man für die Rang-2-Näherung speichern ($k(m+n)$) im Vergleich zu $m\cdot n$?
\end{enumerate}
\karo{8cm}
>>>>>>> fa386f7 (commit)

\end{document}
